% Include this section in a larger article after loading amsmath/amssymb/hyperref.
\section{A verified intermediate with an \texorpdfstring{$h$}{h}-only acceptance test}
\label{sec:h-only-intermediate}
This section records a useful intermediate construction.  Its table still has
196 working states; it is not, by itself, a solution of the strict state-count
problem.  It changes only five rows of the published PA196 table and passes
4,958 uniform source-site contracts for a 3,811-instruction source program.
The construction transfers the $h=1$ acceptance convention already present in
the U208 family into the newer PA196 arithmetic and command language.

\subsection{Command semantics and the change in the source program}
Write a command store as $(t,w,d,h)$.  A nonhalting Shift changes it to
\[
  (h,\operatorname{tail}(w),d,\operatorname{head}(w)),
\]
where head and tail are the two Cantor-decoding components, including the
usual zero-underflow convention.  The published PA196 Shift halts if
$t+h=1$.  The variant halts if $h=1$.  All other commands and all normal
Shift effects are unchanged.

The low-level Shift preparation assumes that its work counter $X$ is zero.
Its original three-instruction prefix transfers $T$ into $X$.  Replace that
prefix by the two-instruction loop clearing $T$.  The rest of the preparation
copies $H$ into $T$ and $X$, transfers $X$ back to $H$, increments $X$, and
returns to the common comparison routine.  Consequently its register effect
before the comparison is
\[
\begin{array}{c|ccc}
 &T&H&X\\\hline
\text{original}&h&t+h&1\\
\text{variant}&h&h&1.
\end{array}
\]
The comparison halts exactly when $H=X$.  On its normal exit it clears the
comparison counters and proceeds to the same stack pop.  This proves the
claimed command semantics by the existing elementary transfer and comparison
lemmas; no claim about arithmetic consistency enters this calculation.

In original source numbering the transfer loop occupies sites 3800--3802.
The variant uses site 3800, $\operatorname{dec}T$, followed by site 3801,
$\operatorname{continueTo}(1)$.  The $\operatorname{inc}X$ is removed and the
loop-back depth changes from two to one.  All old sites after 3802 have new
index one smaller; all sites through 3799 keep their ordinal index.  The
resulting ordinal control graph agrees with this explicit edit.

\subsection{The necessary reassociation}
A naive call to the existing clear-$T$ decoder does not work: that decoder
shares a physical state with clear-$H$ and relies on the preceding address
bit.  Merely substituting the source subtree therefore fails two source-site
contracts.  A first reassociation fixes clear-$T$ but breaks the preceding
\texttt{continueTo(3)}.  The following wider reassociation fixes both contexts.

Let $A=\texttt{node36}$, $B=\texttt{node167}$,
$C=\texttt{node145}$ (clear $T$), and $D=\texttt{node57}$.
Replace the old subtree
\[
  (A,(B,(\texttt{node37},D)))
\quad\text{by}\quad
  ((A,(B,C)),D).
\]
Here $\texttt{node37}$ is the old $T$-to-$X$ transfer.  No control-transfer
operand elsewhere in the tree changes.  With the published dense table
labels, the complete row patch is
\[
\begin{array}{cc|ccc}
\text{state}&\text{read}&\text{write}&\text{move}&\text{next}\\\hline
37&0&0&R&11\\
37&1&1&R&101\\
127&0&0&R&133\\
127&1&1&R&20\\
133&0&0&R&120.
\end{array}
\]
Every unlisted transition is exactly the PA196 transition.  The original
halting state remains 196, the initial state remains 0, and the physical
source-root state remains 75.  All 196 working states remain graph reachable
from state 0.  The maximum source address length grows from 38 to 39, still
within the unchanged 48-cell header.

\subsection{Finite and uniform verification}
The file \texttt{h\_catalan\_node170\_3\_133\_37.tm} contains the complete table;
the adjacent \texttt{\_sites.json} file gives the ordered source sites.
The symbolic runtime check covers every increment, every jump, every halt,
and both zero and positive branches of every decrement.  It verifies 4,958
contracts.  The check treats the padding length as an arbitrary natural
parameter and the entire register file as an opaque symbolic block.  It
therefore does not bound the input values, the register sizes, or the header
padding.  The register-block summaries use the same protected engine rows as
PA196; those rows have not changed.

An independent sparse-tape interpreter, using only literal binary steps,
finds the exact canonical initial source frame after 57,922 steps from the
blank tape.  It checks the complete set of nonzero tape cells, not just a
finite window.  The head is at absolute position $-2$, and the current state
is 75.  Its control address is $1000000$, whose extra final zero is absorbed
by the zero padding.  The physical register file is
\[
 [0,0,0,0,0,0,2,0,0]\mathbin{+\!+}[1]^{26}.
\]
The same interpreter independently reproduces the published PA196 startup
length of 58,154 steps.  These receipts are in
\texttt{h\_symbolic\_result.json} and \texttt{h\_startup\_result.json}.
They are reproducible proof calculations using the published runtime lemmas,
not a claim that a newly generated Lean proof has been kernel checked.

\subsection{Preservation of the consistency property}
The PA196 proof invariant is syntactic: $t$, $h$, and every stack member have
finite derivations in its numerical calculus.  It is not a truth predicate.
The invariant is preserved by the new normal Shift for precisely the same
reason as by the old one.  A new halt has $h=1$, hence supplies a finite
derivation of the distinguished contradiction code.  The syntactic proof
translation into PA described separately converts this into a PA refutation.
Every operation on a bounded execution and its attached derivations is
primitive recursive, and its correctness can be proved in PA.  This yields
\[
  \mathrm{PA}\vdash
  \operatorname{Halt}(M_h)\longrightarrow\neg\operatorname{Con}(\mathrm{PA}).
\]

For the other direction, start with the existing compiler from a PA
refutation to an old halting command word.  Take the execution prefix
\emph{before} its first halting Shift.  Replay that prefix with the new
semantics.  If a new Shift accepts sooner, the desired halt is already
obtained.  Otherwise the two stores agree throughout the prefix.  At its
end, $t+h=1$, so exactly one of the following cases occurs.
\begin{enumerate}
\item If $h=1$, append one Shift, which accepts under the new convention.
\item If $t=1$ and $h=0$, discard the old final Shift and append
      $\mathrm{AX}(0,0,0)$ followed by ten Shifts.
\end{enumerate}
In the second case the AX command has false gate, sets $t=h=0$, and pushes
its eight generated heads followed by the saved old top $1$.  All eight
heads differ from 1: seven have outer implication tag 0 and the remaining
head has universal tag 2.  After nine normal Shifts the new $h$ equals the
saved 1; the tenth Shift accepts.  The older stack contents do not matter.
Applying the AX command \emph{after} replaying the old final Shift would be
incorrect because that Shift can discard the saved $t=1$; the construction
explicitly stops before it.

This padding operation is primitive recursive in the finite old execution.
Together with exhaustive command-word enumeration and the positive finite
runtime simulation, it proves within PA
\[
  \neg\operatorname{Con}(\mathrm{PA})
       \longrightarrow\operatorname{Halt}(M_h).
\]
Thus, at the ordinary mathematical proof level, the intermediate has the
same full PA-provable consistency equivalence as the target construction.
The statement uses the separately supplied syntactic converse argument;
it must not be attributed to the unmodified upstream Lean theorem, which
explicitly states only its existing direction.
